%% ma: L12-B12-0006
%% lop: 12
%% chuong: 4
%% bai: 12
%% dang: 12.12.3
%% loai: TN4
%% muc: NB
%% dap_an: "A"
%% nguon: "(NBV)-ĐỀ SỐ 3-MỖI NGÀY 1 ĐỀ THI 2025.pdf (D03, MỖI NGÀY 1 ĐỀ - Đề số 3), Phần I Câu 6"
%% kien_thuc: "Bài 12 (tích phân, công thức Newton–Leibniz)"
%% trang_thai: chinh-thuc
%% ghi_chu: "Đáp án gốc A đúng. Giáo viên duyệt gộp dạng 12.12.4 vào 12.12.3 (10/10/2026)"
\begin{ex}
Cho $F(x)$ là một nguyên hàm của hàm số $f(x)$ trên $\mathbb R$ thoả mãn $F(1)=-1$, $F(3)=21$. Giá trị của $\displaystyle\int_1^3f(x)\,\mathrm dx$ bằng
\choice{\True $22$}{$20$}{$-21$}{$21$}
\loigiai{$\displaystyle\int_1^3f(x)\,\mathrm dx=F(3)-F(1)=21-(-1)=22$.}
\end{ex}
